\section{Опыт работы}
\resumeSubHeadingListStart
\resumeSubheadingCompany
{\href{https://rsreu.ru/vuz/structure/nii-foton}{\underline{НИИ обработки аэрокосмических изображений <<Фотон>>}}}{Июль 2023 -- Июль 2026}

\resumeSubrole{\currentRole}
\small{Отвечал за полный цикл создания ML-компонента системы обнаружения изменений на аэрокосмических снимках площадной съемки:
	сформировал и разметил обучающий датасет по данным космических аппаратов, обучил нейросетевую модель сегментации
	UNet\plusplus~на ResNet-34 с помощью PyTorch.
	Разработал прикладное ПО на \cpp~(с GUI на Qt), интегрировав в него обученную модель в формате ONNX и
	объединив детерминированные алгоритмы компьютерного зрения с методами глубокого обучения.}

\vspace{4pt}

\resumeSubrole{R\&D-инженер}
\small{Проводил научно-исследовательские работы в области обработки аэрокосмических изображений.
	Устранил появление артефактов при сшивке мозаичных изображений сложных сцен,
	разработав алгоритм формирования стоимостной карты штрафов для корректного проведения линии пореза.
	Выполнил анализ предметной области, по материалам исследований подготовил и опубликовал 4 научные статьи.}
\resumeSubHeadingListEnd
